\documentclass[12pt,a4paper]{article}

% --- Paquetes de Idioma y Formato ---
\usepackage[utf8]{inputenc}
\usepackage[spanish]{babel}
\usepackage[margin=2.5cm]{geometry}
\usepackage{xcolor}
\usepackage{listings}
\usepackage{amsmath}

% --- Configuración de colores para el código MATLAB ---
\definecolor{codegreen}{rgb}{0,0.6,0}
\definecolor{codegray}{rgb}{0.5,0.5,0.5}
\definecolor{codepurple}{rgb}{0.58,0,0.82}
\definecolor{backcolour}{rgb}{0.97,0.97,0.95}

\lstset{
    backgroundcolor=\color{backcolour},   
    commentstyle=\color{codegreen},
    keywordstyle=\color{blue},
    numberstyle=\tiny\color{codegray},
    stringstyle=\color{codepurple},
    basicstyle=\ttfamily\footnotesize,
    breakatwhitespace=false,         
    breaklines=true,                 
    captionpos=b,                    
    keepspaces=true,                 
    numbers=left,                    
    numbersep=5pt,                  
    showspaces=false,                
    showstringspaces=false,
    showtabs=false,                  
    tabsize=2,
    language=Matlab,
    frame=single
}

\begin{document}

% --- Portada ---
\begin{titlepage}
    \centering
    \vspace*{2cm}
    {\Huge \textbf{Informe de Programación}} \\
    \vspace{0.5cm}
    {\Large Algoritmo de Verificación de Números Primos} \\
    \vspace{2cm}
    \textbf{Materia:} [Nombre de la Materia] \\
    \textbf{Autor:} [Tu Nombre] \\
    \vfill
    \textbf{Fecha:} \today
\end{titlepage}

\section{Introducción}
El presente documento detalla la implementación de un script en el entorno MATLAB diseñado para identificar números primos y generar secuencias numéricas descendentes. El objetivo es proporcionar una herramienta sencilla pero robusta para el análisis básico de números naturales.

\section{Metodología y Funciones}
Para la detección de números primos, se ha utilizado la función integrada de MATLAB \texttt{isprime(n)}. Esta función utiliza algoritmos de primalidad que retornan un valor booleano ($1$ para verdadero, $0$ para falso). 

Adicionalmente, se emplea el operador de dos puntos (\texttt{:}) para la creación de vectores decrementales, permitiendo listar los números anteriores al valor ingresado de forma eficiente.

\section{Código Fuente}
A continuación se presenta la implementación final del script:

\begin{lstlisting}[language=Matlab, caption=Script de detección de primos y cuenta regresiva]
% 1. Solicitar el número al usuario
n = input('Introduce un número entero: ');

% 2. Verificación de primalidad
if isprime(n)
    fprintf('El número %d es PRIMO.\n', n);
else
    fprintf('El número %d NO es primo.\n', n);
end

% 3. Lógica para mostrar números anteriores
if n > 1
    fprintf('Números anteriores: ');
    % Creamos el vector descendente de n-1 hasta 1
    anteriores = (n-1):-1:1; 
    fprintf('%d ', anteriores);
    fprintf('\n');
else
    fprintf('No hay números naturales anteriores a %d.\n', n);
end
\end{lstlisting}

\section{Recomendaciones de Uso}
Para garantizar el correcto funcionamiento y la legibilidad del código, se sugieren las siguientes prácticas:
\begin{itemize}
    \item \textbf{Validación:} Asegurarse de que el valor ingresado sea un número entero positivo para evitar errores en la función \texttt{isprime}.
    \item \textbf{Documentación:} Mantener los comentarios en el código para facilitar el mantenimiento por parte de otros usuarios.
    \item \textbf{Formato en LaTeX:} Utilizar el paquete \texttt{listings} (como se muestra en este informe) para que el código sea legible y profesional en documentos impresos.
\end{itemize}

\section{Conclusión}
El programa cumple con los requisitos planteados, demostrando la versatilidad de MATLAB para el manejo de lógica condicional y estructuras de datos simples como los vectores.

\end{document}